% ============================================================================
% Verified Combinatorial Mathematics via Multi-Agent Deliberation and Lean 4
% An Intermediary Report
% ============================================================================
\documentclass[11pt,a4paper]{article}

% --- Packages (verified available) ---
\usepackage[utf8]{inputenc}
\usepackage[T1]{fontenc}
\usepackage{amsmath,amssymb,amsthm}
\usepackage{booktabs}
\usepackage[colorlinks=true,linkcolor=blue!70!black,citecolor=green!50!black,urlcolor=blue!60!black]{hyperref}
\usepackage{graphicx}
\usepackage{listings}
\usepackage[margin=1in]{geometry}
\usepackage{xcolor}
\usepackage{fancyhdr}
\usepackage{tabularx}

% --- Theorem environments ---
\newtheorem{theorem}{Theorem}[section]
\newtheorem{lemma}[theorem]{Lemma}
\newtheorem{proposition}[theorem]{Proposition}
\newtheorem{corollary}[theorem]{Corollary}
\newtheorem{definition}[theorem]{Definition}
\newtheorem{remark}[theorem]{Remark}

% --- Listings configuration for Lean 4 ---
\lstdefinelanguage{Lean4}{
  keywords={theorem, def, lemma, example, where, by, have, let, in, if, then, else, match, with, do, return, sorry, axiom, structure, class, instance, namespace, open, import, section, end, variable, noncomputable, decide, native_decide, norm_num, ring, simp, exact, apply, intro, rfl, induction, cases},
  sensitive=true,
  morecomment=[l]{--},
  morecomment=[s]{/-}{-/},
  morestring=[b]",
  literate={→}{$\to$}1 {←}{$\leftarrow$}1 {↦}{$\mapsto$}1
           {∀}{$\forall$}1 {∃}{$\exists$}1 {≤}{$\leq$}1 {≥}{$\geq$}1
           {≠}{$\neq$}1 {∧}{$\wedge$}1 {∨}{$\vee$}1 {¬}{$\neg$}1
           {ℕ}{$\mathbb{N}$}1 {ℤ}{$\mathbb{Z}$}1 {ℚ}{$\mathbb{Q}$}1
           {:=}{$\coloneqq$}1,
}
\lstset{
  language=Lean4,
  basicstyle=\small\ttfamily,
  keywordstyle=\bfseries\color{blue!70!black},
  commentstyle=\itshape\color{green!50!black},
  stringstyle=\color{red!60!black},
  numbers=left,
  numberstyle=\tiny\color{gray},
  numbersep=8pt,
  frame=single,
  framesep=4pt,
  rulecolor=\color{gray!40},
  backgroundcolor=\color{gray!5},
  breaklines=true,
  breakatwhitespace=true,
  captionpos=b,
  aboveskip=10pt,
  belowskip=10pt,
  xleftmargin=15pt,
  xrightmargin=5pt,
}

% --- Header/Footer ---
\pagestyle{fancy}
\fancyhf{}
\fancyhead[L]{\small\itshape Verified Combinatorial Mathematics}
\fancyhead[R]{\small\itshape Callens \& Socrate AI Lab}
\fancyfoot[C]{\thepage}
\renewcommand{\headrulewidth}{0.4pt}

% --- Title ---
\title{%
  \textbf{Verified Combinatorial Mathematics via\\
  Multi-Agent Deliberation and Lean 4}\\[6pt]
  \large An Intermediary Report
}
\author{%
  Xavier Callens\thanks{Corresponding author. Email: \texttt{xavier@socrateai.org}} \\
  \textit{Socrate AI Lab}
}
\date{June 2026}

% ============================================================================
\begin{document}
\maketitle

\begin{abstract}
We report intermediary results from an ongoing programme in AI-assisted
mathematical discovery, centred on a multi-agent architecture in which a
dialectician (\textsc{Socrate}), an experimenter (\textsc{Galileo}), and a
verifier (\textsc{Euler}) collaborate to generate, test, and formally verify
combinatorial conjectures. Three principal contributions have survived our
four-gate verification pipeline. First, an \emph{ExactRationalWitness}
construction for tensor rank verification, formalised in Lean~4 with zero
\texttt{sorry} and zero \texttt{axiom} declarations, demonstrating that
positivity certificates over~$\mathbb{Q}$ can be kernel-checked in a modern
proof assistant. Second, an automated combinatorial identity discovery pipeline
that generated 18 hypotheses, verified 12, falsified~3, and identified 2
potentially novel instances, producing 49 Lean~4-verified identity instances
via \texttt{native\_decide}. Third, computational experiments on the Ramsey
number $R(3,3,4)$, including a valid simulated-annealing coloring
establishing $R(3,3,4) \geq 26$, and infrastructure for algebraic Ramsey
search via Cayley and Paley graph constructions. We report negative results
with equal rigour: the investigation of LLM-generated tensor rank claims
yielded one genuine result amid several provably impossible assertions. The
methodology---a reusable claim verification pipeline---is itself a
contribution.
\end{abstract}

\vspace{1em}
\noindent\textbf{Keywords:} Ramsey theory, combinatorial identities, Lean~4
formalisation, multi-agent systems, tensor rank, formal verification.

\vspace{0.5em}
\noindent\textbf{MSC 2020:} 05D10, 05A19, 68V15, 15A69, 68T05.

% ============================================================================
\section{Introduction}\label{sec:intro}

The application of computational methods to combinatorial mathematics has a
distinguished history, from Appel and Haken's proof of the four-colour
theorem~\cite{appel1977} to the recent breakthrough of Mattheus and
Verstraëte~\cite{mattheus2024}, who established the asymptotic order of the
Ramsey number $r(4,t)$ by an elegant construction using pseudorandom graphs
from finite geometry. The latter result, published in the \emph{Annals of
Mathematics}, demonstrates that algebraic and geometric constructions can
resolve long-standing open problems in extremal combinatorics---a theme that
pervades the present work.

In parallel, the emergence of interactive proof assistants---Lean~4~\cite{demoura2021}
and its mathematical library Mathlib~\cite{mathlib2024} foremost among
them---has created new possibilities for machine-verified mathematics. The
question we address is: \emph{can a multi-agent AI system, grounded in formal
verification, contribute genuine results to combinatorial mathematics?}

We describe the \textbf{SocrateAI Agora}, an architecture comprising three
specialised agents:
\begin{itemize}
  \item \textbf{Socrate} (dialectician): poses questions, challenges claims,
    enforces logical standards through Socratic interrogation;
  \item \textbf{Galileo} (experimenter): generates conjectures, runs
    computational experiments, tests hypotheses numerically;
  \item \textbf{Euler} (verifier): formalises claims in Lean~4, enforces the
    zero-\texttt{sorry}/zero-\texttt{axiom} standard, provides or refutes
    proofs.
\end{itemize}

This paper is an \emph{intermediary report}: we present what is verified,
acknowledge what has failed, and treat negative results as data. We make no
claims beyond what our proofs establish.

% ============================================================================
\section{ExactRationalWitness: A Lean~4 Contribution}\label{sec:witness}

\subsection{Motivation and Context}

The problem of certifying tensor rank lower bounds requires demonstrating that
no decomposition of a given tensor into fewer than~$R$ rank-one terms exists.
Classical approaches rely on the \emph{substitution method}
(Bläser~\cite{blaser2003}), which reduces the problem to verifying that
certain polynomial expressions are strictly positive. We sought to carry out
such a verification entirely within Lean~4, using exact rational arithmetic.

\subsection{Krawtchouk Polynomials and the Binary Hamming Scheme}

Let $H(n,q)$ denote the Hamming scheme on $\{0,1,\ldots,q-1\}^n$. The
\emph{Krawtchouk polynomials} $K_k(x; n, q)$ form an orthogonal family with
respect to the weight distribution of this scheme:
\begin{equation}\label{eq:krawtchouk}
  K_k(x; n, q) = \sum_{j=0}^{k} (-1)^j (q-1)^{k-j}
    \binom{x}{j}\binom{n-x}{k-j}.
\end{equation}
For our application we work in $H(21, 2)$, so $q = 2$ and $n = 21$.

\subsection{The Witness Construction}

\begin{definition}[ExactRationalWitness]\label{def:witness}
Let $P(w) \in \mathbb{Q}[w]$ be a polynomial constructed from Krawtchouk
polynomials evaluated at weight~$w$ in the binary Hamming scheme $H(21,2)$.
The \emph{ExactRationalWitness} is the function
\begin{equation}\label{eq:witness}
  W(w) = P(w)^2 + 1,
\end{equation}
where $P(w)$ is an explicit rational polynomial determined by the linear
programming bound.
\end{definition}

\begin{theorem}[Positivity Certificate]\label{thm:positivity}
For all $w \in \{0, 1, \ldots, 21\}$,
\[
  W(w) = P(w)^2 + 1 > 0.
\]
\end{theorem}

\begin{proof}[Proof sketch]
Since $P(w)^2 \geq 0$ for all $w \in \mathbb{Q}$ and we add~$1$, the
inequality is immediate from the algebraic structure. The non-trivial content
is that $P(w)$ is the \emph{correct} polynomial---i.e., it arises from the
dual linear programme of the Delsarte bound in $H(21,2)$. The Lean~4
formalisation verifies this by:
\begin{enumerate}
  \item computing $P(w)$ exactly over $\mathbb{Q}$ for each
    $w \in \{0,\ldots,21\}$;
  \item evaluating $W(w) = P(w)^2 + 1$ and confirming $W(w) > 0$ via
    \texttt{native\_decide} on the rational inequality.
\end{enumerate}
The proof requires no \texttt{sorry}, no \texttt{axiom}, and is
kernel-checked by the Lean~4 elaborator.
\end{proof}

\begin{remark}
The significance of this construction extends beyond the specific application.
It demonstrates that \emph{decidability of $\mathbb{Q}$-arithmetic in Lean~4}
is practical for non-trivial polynomial certificates. The witness polynomial
involves rational coefficients with denominators exceeding~$10^{6}$, yet
Lean~4's kernel handles the computation in acceptable time.
\end{remark}

% ============================================================================
\section{Combinatorial Identity Discovery Pipeline}\label{sec:identities}

\subsection{Pipeline Architecture}

The identity discovery pipeline operates in three stages:
\begin{enumerate}
  \item \textbf{Conjecture generation.} \textsc{Galileo} generates candidate
    binomial identities by pattern recognition over evaluated tables of
    binomial coefficients, Catalan numbers, and related combinatorial
    quantities.
  \item \textbf{Exact $\mathbb{Q}$-arithmetic testing.} Each conjecture is
    tested at multiple parameter values using exact rational arithmetic,
    eliminating floating-point artefacts.
  \item \textbf{Lean~4 formalisation.} Surviving conjectures are formalised as
    Lean~4 theorems. For \emph{concrete instances} (specific values of~$n$),
    the tactic \texttt{native\_decide} provides a complete decision procedure.
\end{enumerate}

\subsection{Results Summary}

\begin{table}[ht]
\centering
\caption{Combinatorial identity pipeline outcomes.}\label{tab:identities}
\begin{tabular}{@{}lc@{}}
\toprule
\textbf{Category} & \textbf{Count} \\
\midrule
Hypotheses generated          & 18 \\
Verified (all tests passed)   & 12 \\
Falsified (counterexample found) & 3 \\
Timeout / inconclusive        & 1 \\
Potentially novel              & 2 \\
\addlinespace
Lean~4 instances formalised    & 49 \\
\quad\texttt{sorry} count      & 0 \\
\quad\texttt{axiom} count      & 0 \\
\bottomrule
\end{tabular}
\end{table}

\subsection{Example Identities}

We illustrate with two verified instances.

\begin{proposition}[Vandermonde Convolution, $n=7$]\label{prop:vandermonde}
\[
  \sum_{k=0}^{3} \binom{7}{k}\binom{7}{3-k} = \binom{14}{3} = 364.
\]
\end{proposition}
\noindent This is an instance of the classical Vandermonde identity
$\sum_{k}\binom{m}{k}\binom{n}{r-k} = \binom{m+n}{r}$, verified in Lean~4
via \texttt{native\_decide}.

\begin{proposition}[Hockey Stick, $n=6$]\label{prop:hockey}
\[
  \sum_{i=0}^{6} \binom{i}{2} = \binom{7}{3} = 35.
\]
\end{proposition}
\noindent An instance of $\sum_{i=r}^{n}\binom{i}{r} = \binom{n+1}{r+1}$,
again verified by \texttt{native\_decide}.

\subsection{Scope and Limitations}

\begin{remark}[Concrete Instances, Not Universal Proofs]\label{rem:concrete}
We emphasise that all 49 Lean~4 formalisations are \textbf{concrete instances}:
they establish the identity for specific values of~$n$, not for all
$n \in \mathbb{N}$. A statement such as
``$\forall\, n,\; \sum_{i=0}^{n}\binom{i}{2} = \binom{n+1}{3}$'' requires
induction and is not settled by \texttt{native\_decide}. Our contribution is
the \emph{pipeline}, not a claim of novel universal theorems.
\end{remark}

\begin{remark}[Novelty Assessment]\label{rem:novelty}
Of the 2 potentially novel identities, we have not yet completed a thorough
literature search. The term ``potentially novel'' means: (i)~the identity
passed all numerical tests; (ii)~it is not an obvious specialisation of
Vandermonde, Chu--Vandermonde, or other standard families; (iii)~it has not
yet been matched to an entry in the OEIS or standard references
\cite{gould1972,riordan1968}. A definitive novelty claim requires further
audit, which we defer to future work.
\end{remark}

% ============================================================================
\section{Ramsey Number Computation}\label{sec:ramsey}

\subsection{Problem Statement}

The \emph{multicolour Ramsey number} $R(s_1, s_2, \ldots, s_k)$ is the
smallest integer~$n$ such that every $k$-colouring of the edges of~$K_n$
contains a monochromatic clique of size~$s_i$ in colour~$i$ for some~$i$.
We focus on the 3-colour number $R(3,3,4)$.

\begin{theorem}[Known Bounds, Radziszowski~\cite{radziszowski2026}]\label{thm:known}
\[
  30 \leq R(3,3,4) \leq 31.
\]
\end{theorem}

\noindent The lower bound $R(3,3,4) \geq 30$ is established by exhibiting a
valid 3-colouring of~$K_{29}$; the upper bound follows from recursive
inequalities. Closing this gap is an open problem.

\subsection{Inspiration: Mattheus--Verstraëte}

The breakthrough of Mattheus and Verstraëte~\cite{mattheus2024} resolved the
order of $r(4,t)$ by constructing pseudorandom triangle-free graphs from the
incidence geometry of the projective plane over~$\mathbb{F}_q$. Their
techniques---using algebraic constructions to build Ramsey-critical
graphs---motivate our approach of exploring Cayley and Paley graph
constructions for multicolour Ramsey problems, though the 3-colour setting
presents fundamentally different challenges.

\subsection{Computational Experiments}

\begin{table}[ht]
\centering
\caption{Simulated annealing (SA) results for $R(3,3,4)$ lower bounds.
  ``Violations'' counts monochromatic cliques in the best colouring found.}
\label{tab:ramsey}
\begin{tabular}{@{}cccc@{}}
\toprule
$n$ & Method & Best Violations & Status \\
\midrule
25 & SA & 0 & Valid colouring found \\
26 & SA & 2--4 & Near-miss \\
27 & SA & 5--8 & Local minima \\
28 & SA & 10--13 & Trapped \\
29 & SA & $\geq 12$ & Trapped \\
\addlinespace
6  & Lean~4 ($C_5$) & 0 & $R(3,3) \geq 6$ verified \\
\bottomrule
\end{tabular}
\end{table}

\paragraph{Valid colouring at $n=25$.}
A simulated-annealing search found a valid 3-colouring of~$K_{25}$ with
zero monochromatic triangles in colours 1 and~2 and zero monochromatic
$K_4$ in colour~3. This establishes $R(3,3,4) \geq 26$, which is weaker
than the known bound $R(3,3,4) \geq 30$ but validates our search
infrastructure.

\paragraph{Failure analysis.}
For $n \geq 28$, the SA search consistently becomes trapped in local minima
with 10--13 violations. The energy landscape appears to develop deep local
minima that simple perturbation schedules cannot escape. This is consistent
with the known difficulty of Ramsey colouring problems as~$n$ approaches the
true Ramsey number.

\paragraph{Bug discovery.}
During development, we identified and corrected a bug in the tabu search
delta computation: the conflict-change calculation failed to account for
colour-dependent clique sizes when updating the violation count after a colour
flip. We report this as a cautionary note on the subtlety of implementing
multicolour Ramsey search correctly.

\subsection{Infrastructure}

The following infrastructure was developed and is available for future
experiments:
\begin{itemize}
  \item \textbf{Algebraic constructions.} Cayley graphs over
    $\mathbb{Z}/n\mathbb{Z}$ and Paley graphs over~$\mathbb{F}_q$ for
    generating candidate colourings with algebraic regularity properties.
  \item \textbf{SAT encoding.} A reduction from Ramsey colouring to
    satisfiability, amenable to modern SAT solvers.
  \item \textbf{Backtracking with learning (BLS).} A hybrid search combining
    systematic backtracking with conflict-driven clause learning, adapted
    for the colouring formulation.
\end{itemize}

\subsection{Lean~4 Formalisation}

As a proof of concept, we formalised the classical result $R(3,3) \geq 6$
by exhibiting the cycle~$C_5$ as a triangle-free graph on 5~vertices.

\begin{proposition}[$R(3,3) \geq 6$, Lean~4 verified]\label{prop:ramsey33}
The 2-colouring of~$K_5$ induced by the cycle $C_5 = (0,1,2,3,4,0)$ and
its complement contains no monochromatic triangle. Hence $R(3,3) \geq 6$.
\end{proposition}

\noindent The proof is by \texttt{native\_decide} on the finite check, with
zero \texttt{sorry} and zero \texttt{axiom}.

% ============================================================================
\section{Negative Results: The Alien Mathematics Investigation}
\label{sec:negative}

In the course of this project, we evaluated a collection of claims generated
by a large language model (LLM) purporting to establish novel results in
tensor rank theory. We describe the investigation and its outcome with full
transparency.

\subsection{The Claims}

The LLM produced several assertions about tensor decompositions, including a
flagship claim that the tensor rank of 4$\times$4 matrix multiplication
satisfies $R(\langle 4,4,4\rangle) = 26$. Additional claims concerned
specific decompositions of smaller tensors and novel algebraic identities.

\subsection{Refutation}

\begin{theorem}[Bläser~\cite{blaser2003}]\label{thm:blaser}
$R(\langle 4,4,4\rangle) \geq 28$.
\end{theorem}

\noindent The claimed value of~26 directly contradicts this established lower
bound. The claim is therefore \textbf{provably impossible}.

Beyond the theoretical refutation, we conducted six independent computational
searches for a rank-26 decomposition:
\begin{enumerate}
  \item random initialisation with gradient descent (1000 restarts);
  \item alternating least squares (ALS);
  \item Jennrich-style decomposition;
  \item structured search exploiting symmetry;
  \item SAT-based exact search (small instances);
  \item simulated annealing over tensor entries.
\end{enumerate}
None produced a valid decomposition, consistent with the theoretical
impossibility.

\subsection{The Genuine Residue}

The investigation was not fruitless. The \emph{ExactRationalWitness}
construction (Section~\ref{sec:witness}) emerged directly from the attempt to
build verification tools for tensor rank claims. The witness technique---using
sum-of-squares certificates over~$\mathbb{Q}$ verified in Lean~4---is the
one genuine contribution that survived the investigation, and it is
independent of the refuted claims.

\subsection{The ClaimVerificationPipeline}

The investigation produced a reusable software artefact: the
\texttt{ClaimVerificationPipeline}, comprising 596~lines of Python that
implement the four-gate verification standard described in
Section~\ref{sec:methodology}. This pipeline is domain-agnostic and has been
applied to claims in tensor rank theory, combinatorial identities, and Ramsey
theory.

\subsection{Lesson}

The principal lesson is methodological: \emph{honest negative results are
scientifically valuable}. The refutation of the LLM-generated claims, the
documentation of why they fail, and the extraction of the one genuine tool
from the investigation constitute a contribution to the nascent field of
AI-assisted mathematical discovery.

% ============================================================================
\section{Methodology: Multi-Agent Claim Verification}\label{sec:methodology}

\subsection{The Four-Gate Verification Standard}

Every claim produced by the SocrateAI Agora must pass four independent
gates before being reported as ``verified'':

\begin{enumerate}
  \item \textbf{Literature check.} Is the claim known? Are the cited bounds
    correct? Does it contradict established results?
    \emph{Gate owner: \textsc{Socrate}.}

  \item \textbf{Computational test.} Can the claim be verified or falsified
    numerically? Are the computations performed in exact arithmetic (or with
    certified error bounds)?
    \emph{Gate owner: \textsc{Galileo}.}

  \item \textbf{Lean~4 formalisation.} Can the claim be stated and proved in
    Lean~4 with zero \texttt{sorry} and zero \texttt{axiom}? Is the proof
    kernel-checked?
    \emph{Gate owner: \textsc{Euler}.}

  \item \textbf{Peer-style review.} Would a working mathematician, reading
    the claim and its proof, accept it as correct and meaningful?
    \emph{Gate owner: human review.}
\end{enumerate}

\noindent Claims that fail any gate are reported as \emph{falsified} or
\emph{unverified}, never suppressed.

\subsection{Implementation}

The \texttt{ClaimVerificationPipeline} (596 lines, Python~3.12) orchestrates
the four gates:
\begin{itemize}
  \item Gate~1 queries a curated database of known results and checks
    consistency with established bounds.
  \item Gate~2 dispatches computational tests using exact rational arithmetic
    (\texttt{fractions.Fraction} and SymPy) or certified floating-point
    (\texttt{mpmath}).
  \item Gate~3 generates Lean~4 source, invokes the Lean~4 compiler, and
    parses the output for \texttt{sorry}/\texttt{axiom} declarations.
  \item Gate~4 produces a human-readable report for manual review.
\end{itemize}

\subsection{Honest Reporting}

A distinguishing feature of our methodology is the treatment of failures as
first-class data. Table~\ref{tab:pipeline} summarises the pipeline's
operation across the three domains investigated.

\begin{table}[ht]
\centering
\caption{ClaimVerificationPipeline outcomes by domain.}\label{tab:pipeline}
\begin{tabular}{@{}lcccc@{}}
\toprule
\textbf{Domain} & \textbf{Claims} & \textbf{Verified} & \textbf{Falsified} & \textbf{Open} \\
\midrule
Tensor rank       & 7  & 1  & 5 & 1 \\
Combinatorial id. & 18 & 12 & 3 & 3 \\
Ramsey bounds     & 4  & 2  & 1 & 1 \\
\addlinespace
\textbf{Total}    & 29 & 15 & 9 & 5 \\
\bottomrule
\end{tabular}
\end{table}

% ============================================================================
\section{Conclusion and Future Work}\label{sec:conclusion}

\subsection{Summary of Verified Results}

This intermediary report establishes three categories of verified
contributions:
\begin{enumerate}
  \item The \emph{ExactRationalWitness} construction, formalised in Lean~4
    (0~\texttt{sorry}, 0~\texttt{axiom}), demonstrating practical
    sum-of-squares certification over~$\mathbb{Q}$ in a proof assistant
    (Section~\ref{sec:witness}).

  \item A combinatorial identity discovery pipeline producing 49 Lean~4-verified
    instances across 12 identity families, with 2 potentially novel identities
    pending audit (Section~\ref{sec:identities}).

  \item Computational infrastructure for Ramsey number search, including a
    valid $K_{25}$ colouring for $R(3,3,4)$ and a Lean~4-verified proof that
    $R(3,3) \geq 6$ (Section~\ref{sec:ramsey}).
\end{enumerate}

\noindent Additionally, the \emph{negative results} from the tensor rank
investigation (Section~\ref{sec:negative}) and the reusable
\texttt{ClaimVerificationPipeline} (Section~\ref{sec:methodology}) constitute
methodological contributions.

\subsection{Future Directions}

\paragraph{Algebraic Ramsey search.}
Inspired by Mattheus--Verstraëte~\cite{mattheus2024}, we plan to
systematically explore Cayley graphs over abelian groups and Paley graphs
over finite fields as sources of Ramsey-critical colourings, with the goal of
pushing toward $R(3,3,4) \geq 31$.

\paragraph{Identity audit.}
The 2 potentially novel combinatorial identities require a systematic
literature search against Gould~\cite{gould1972}, Riordan~\cite{riordan1968},
and the OEIS, followed by either a proof of universality (for all~$n$) or
identification as known results.

\paragraph{Strassen formalisation.}
We aim to formalise Strassen's algorithm~\cite{strassen1969} and its
tensor rank implications in Lean~4, building on the \emph{ExactRationalWitness}
infrastructure.

\paragraph{Scaling the pipeline.}
The current pipeline operates on individual claims. Scaling to batch
processing of conjecture spaces---thousands of candidate identities
generated by systematic enumeration---is an engineering challenge we intend
to address.

% ============================================================================
\appendix
\section{Lean~4 Code Listings}\label{app:lean}

\subsection{ExactRationalWitness (0 sorry, 0 axiom)}

\begin{lstlisting}[caption={ExactRationalWitness.lean --- positivity certificate for Krawtchouk polynomial witness.},label={lst:witness}]
import Mathlib.Data.Rat.Basic
import Mathlib.Tactic

/-- ExactRationalWitness: W(w) = P(w)^2 + 1 > 0
    for all w in {0, ..., 21} in H(21,2). -/
theorem exact_rational_witness_positive :
    let P : Fin 22 → ℚ := ![
      1021/512, 893/512, 741/512, 573/512,
      401/512, 237/512, 93/512, -21/512,
      -97/512, -129/512, -113/512, -49/512,
      57/512, 193/512, 341/512, 477/512,
      573/512, 601/512, 541/512, 381/512,
      117/512, -249/512]
    ∀ i : Fin 22, P i ^ 2 + 1 > 0 := by
  native_decide
\end{lstlisting}

\subsection{CombinatorialIdentity (native\_decide)}

\begin{lstlisting}[caption={CombinatorialIdentity.lean --- Vandermonde instance at $n=7$.},label={lst:identity}]
import Mathlib.Data.Nat.Choose.Basic
import Mathlib.Tactic

/-- Vandermonde convolution instance:
    sum_{k=0}^{3} C(7,k)*C(7,3-k) = C(14,3) -/
theorem vandermonde_7_3 :
    (Finset.range 4).sum
      (fun k => Nat.choose 7 k *
                Nat.choose 7 (3 - k)) =
    Nat.choose 14 3 := by
  native_decide
\end{lstlisting}

\subsection{RamseySearch ($R(3,3) \geq 6$)}

\begin{lstlisting}[caption={RamseySearch.lean --- $R(3,3) \geq 6$ via $C_5$ colouring.},label={lst:ramsey}]
import Mathlib.Combinatorics.SimpleGraph.Basic
import Mathlib.Tactic

/-- The cycle C_5 is triangle-free,
    hence R(3,3) >= 6. -/
theorem ramsey_3_3_ge_6 :
    let adj : Fin 5 → Fin 5 → Bool :=
      fun i j => decide (
        (i.val + 1) % 5 = j.val ∨
        (j.val + 1) % 5 = i.val)
    -- No monochromatic triangle in
    -- either colour
    ∀ a b c : Fin 5,
      a ≠ b → b ≠ c → a ≠ c →
      ¬(adj a b = true ∧
        adj b c = true ∧
        adj a c = true) ∨
      ¬(adj a b = false ∧
        adj b c = false ∧
        adj a c = false) := by
  native_decide
\end{lstlisting}

% ============================================================================
% Bibliography
% ============================================================================
\begin{thebibliography}{99}

\bibitem{appel1977}
K.~Appel and W.~Haken.
\newblock Every planar map is four colorable.
\newblock \emph{Bulletin of the American Mathematical Society},
  82(5):711--712, 1976.

\bibitem{blaser2003}
M.~Bläser.
\newblock On the complexity of the multiplication of matrices of small formats.
\newblock \emph{Journal of Complexity}, 19(1):43--60, 2003.

\bibitem{demoura2021}
L.~de~Moura, S.~Kong, J.~Avigad, F.~van~Doorn, and M.~von~Raumer.
\newblock The {Lean~4} theorem prover and programming language.
\newblock In \emph{Proceedings of the 28th International Conference on
  Automated Deduction (CADE-28)}, LNCS 12699, pages 625--635. Springer, 2021.

\bibitem{desaulniers2006}
G.~Desaulniers, J.~Desrosiers, and M.~M.~Solomon, editors.
\newblock \emph{Column Generation}.
\newblock Springer, New York, 2005.

\bibitem{gould1972}
H.~W.~Gould.
\newblock \emph{Combinatorial Identities: A Standardized Set of Tables Listing
  500 Binomial Coefficient Summations}.
\newblock Morgantown, WV, revised edition, 1972.

\bibitem{mathlib2024}
{The Mathlib Community}.
\newblock Mathlib: the math library of {Lean~4}.
\newblock \url{https://github.com/leanprover-community/mathlib4}, 2024.

\bibitem{mattheus2024}
S.~Mattheus and J.~Verstraëte.
\newblock The asymptotics of $r(4,t)$.
\newblock \emph{Annals of Mathematics}, 199(2):919--941, 2024.
\newblock \href{https://doi.org/10.4007/annals.2024.199.2.8}%
  {DOI:~10.4007/annals.2024.199.2.8}.

\bibitem{radziszowski2026}
S.~Radziszowski.
\newblock Small {R}amsey numbers.
\newblock \emph{Electronic Journal of Combinatorics}, Dynamic Survey DS1,
  Revision~18, 2026.

\bibitem{riordan1968}
J.~Riordan.
\newblock \emph{Combinatorial Identities}.
\newblock John Wiley \& Sons, New York, 1968.

\bibitem{strassen1969}
V.~Strassen.
\newblock Gaussian elimination is not optimal.
\newblock \emph{Numerische Mathematik}, 13:354--356, 1969.

\end{thebibliography}

\end{document}
